\documentclass[11pt]{article}
\usepackage[T1]{fontenc}
\usepackage{lmodern,amsmath,amssymb,amsthm,mathtools}
\usepackage[margin=1in]{geometry}
\usepackage{microtype,booktabs,array,xcolor}
\usepackage[colorlinks=true,linkcolor=blue!45!black,citecolor=blue!45!black,urlcolor=blue!45!black]{hyperref}
\hypersetup{pdftitle={A positive counterexample to the premagic diagonalizability conjecture},pdfauthor={Myeongjin Kang}}
\newtheorem{theorem}{Theorem}
\newtheorem{proposition}[theorem]{Proposition}
\newtheorem{lemma}[theorem]{Lemma}
\theoremstyle{remark}\newtheorem{remark}[theorem]{Remark}
\newcommand{\R}{\mathbb R}
\newcommand{\C}{\mathbb C}
\newcommand{\one}{\mathbf 1}
\DeclareMathOperator{\diag}{diag}
\DeclareMathOperator{\spanof}{span}
\title{A positive counterexample to the\\premagic diagonalizability conjecture}
\author{Myeongjin Kang}
\date{October 25, 2026}
\begin{document}
\maketitle
\begin{abstract}
A square matrix is premagic when its vector of row sums equals its
vector of column sums. We give a strictly positive integer matrix of
order three with every row and column sum equal to nine that is not
diagonalizable, even over the complex numbers. This refutes Conjecture~7.2
of Teknomo's \emph{Premagic and Ideal Flow Matrices}, also indexed as
MathDB problem 336488. A scalar shift gives invertible positive
counterexamples; normalization gives strictly positive doubly stochastic
counterexamples. The obstruction is an explicit Jordan chain of length
two. Accompanying Lean proofs certify the counterexample and its
non-diagonalizability without additional axioms or stipulated research
theorems.
\end{abstract}
\noindent\textit{Status and scope.} This is an AI-assisted research note
with internal independent checks, not external peer review. The
conjecture is refuted directly; no worldwide-priority claim is made.
Original prose is licensed CC BY 4.0 and original proof code MIT.

\section{The question and the counterexample}
Let $A=(a_{ij})\in\R^{n\times n}$ and let $\one=(1,\ldots,1)^T$.
Following Teknomo's Definition~4.3~\cite{Teknomo}, call $A$ \emph{premagic}
if
\begin{equation}\label{eq:premagic}
 \sum_{j=1}^n a_{ij}=\sum_{j=1}^n a_{ji}\quad(1\leq i\leq n),
 \qquad\text{equivalently } A\one=A^T\one.
\end{equation}
This does not require all the row sums to be equal to one another.
Our counterexamples satisfy that additional condition as well.
Conjecture~7.2 of the cited paper asserts diagonalizability for premagic
matrices; MathDB~\cite{MathDB} records the same assertion. We use
diagonalizability in its usual sense of similarity to a diagonal matrix.

\begin{theorem}\label{thm:main}
The matrix
\[
 M=\begin{pmatrix}4&4&1\\2&2&5\\3&3&3\end{pmatrix}
\]
is strictly positive, integral and premagic, with all row sums and
column sums equal to $9$. It is not diagonalizable over $\C$, and hence
is not diagonalizable over $\R$.
\end{theorem}

Here is a direct proof, without a characteristic-polynomial calculation.
Set
\[
 u=\begin{pmatrix}1\\-1\\0\end{pmatrix},\qquad
 v=\begin{pmatrix}1\\1\\-2\end{pmatrix}.
\]
Matrix multiplication gives
\begin{equation}\label{eq:chain}
 Mu=0,\qquad Mv=6u\ne0,\qquad M^2v=0.
\end{equation}
For a diagonalizable operator $T$, one always has
$\ker(T^2)=\ker(T)$: this holds for a diagonal matrix, since
$d_i^2x_i=0$ implies $d_ix_i=0$ in a field, and it is invariant under
change of basis. Thus \eqref{eq:chain} excludes diagonalizability over
$\C$. The row sums, column sums and entrywise positivity are immediate
from the displayed matrix. This proves Theorem~\ref{thm:main}.

\section{An invertible family and stochastic normalization}
The failure is not caused by singularity or zero entries.
For real $t\geq0$ define
\begin{equation}\label{eq:family}
 A_t=M+tI_3
 =\begin{pmatrix}t+4&4&1\\2&t+2&5\\3&3&t+3\end{pmatrix},
 \qquad B_t=\frac{A_t}{t+9}.
\end{equation}

\begin{proposition}\label{prop:family}
For every $t\geq0$, $A_t$ is strictly positive and premagic with common
row and column sum $t+9$. The matrix $B_t$ is strictly positive and
doubly stochastic. Neither is diagonalizable over $\C$.
For $t>0$ both are invertible. Their characteristic polynomials are
\[
 \chi_{A_t}(X)=(X-(t+9))(X-t)^2,
 \qquad
 \chi_{B_t}(X)=(X-1)\left(X-\frac{t}{t+9}\right)^2.
\]
\end{proposition}
\begin{proof}
The row and column sums follow from those of $M$, and positivity follows
entry by entry. Division by the positive number $t+9$ gives the
stochastic assertions. Consider the change-of-basis matrix
\[
 P=(\one\ \ u\ \ v)
 =\begin{pmatrix}1&1&1\\1&-1&1\\1&0&-2\end{pmatrix},\qquad \det P=6.
\]
Since $M\one=9\one$, \eqref{eq:chain} yields
\begin{equation}\label{eq:jordan}
 A_tP=P\begin{pmatrix}t+9&0&0\\0&t&6\\0&0&t\end{pmatrix}.
\end{equation}
This proves the characteristic-polynomial formula and
$\det A_t=(t+9)t^2$. The nonzero off-diagonal entry in the repeated
eigenvalue block cannot be removed by similarity: explicitly,
$(A_t-tI)^2v=0$ while $(A_t-tI)v=6u\ne0$.
If $A_t$ were diagonalizable then so would be $A_t-tI$, contradicting
the kernel argument above. Finally, a nonzero scalar multiple of a
matrix is diagonalizable if and only if the matrix is, proving the
claim for $B_t$. Its characteristic polynomial and invertibility also
follow by scaling.
\end{proof}

In particular, an invertible strictly positive integer counterexample
and its doubly stochastic normalization are
\begin{equation}\label{eq:invertible}
 A_1=\begin{pmatrix}5&4&1\\2&3&5\\3&3&4\end{pmatrix},
 \qquad B_1=\frac1{10}\begin{pmatrix}5&4&1\\2&3&5\\3&3&4\end{pmatrix}.
\end{equation}
The eigenvalues of $B_1$ are $1,1/10,1/10$; the eigenspace for $1/10$
has dimension one. The dominant eigenvalue is simple. Thus the
counterexample is fully compatible with the simplicity of the dominant
eigenvalue for positive matrices: the defect occurs at a different
eigenvalue. No appeal to Perron--Frobenius theory is needed in the proof.

\section{Why the construction works}
Write $J=\one\one^T$. Then
\begin{equation}\label{eq:decomp}
 M=3J+uv^T,
 \qquad \one^Tu=\one^Tv=v^Tu=0.
\end{equation}
Consequently $J(uv^T)=(uv^T)J=0$, and $(uv^T)^2=0$, although
$uv^T\ne0$. The constant part controls the row and column sums, while
the nonzero nilpotent part acts on the zero-sum plane. Positivity of
their sum does not remove this nilpotent part.

This also explains the construction beyond dimension three. For each
$n\geq3$, append zeros to $u$ and $v$, use the $n$-dimensional all-ones
vector, and put
\[
 M_n=3\one\one^T+uv^T.
\]
Every entry is at least $1$, every row and column sum is $3n$, and
$M_nu=0$, $M_nv=6u\ne0$. Thus the same argument provides positive
integer counterexamples in every dimension at least three. This
dimension extension is a mathematical corollary of the displayed
identities; the concrete three-dimensional result is the formal
certificate used to refute the conjecture.

For real matrices dimension three is minimal. In dimension two, the
premagic condition on $\left(\begin{smallmatrix}a&b\\c&d\end{smallmatrix}\right)$
implies $a+b=a+c$, so $b=c$. The matrix is then real symmetric and
diagonalizable by the real spectral theorem. Dimension one is immediate.
This minimality assertion concerns real matrices, not arbitrary complex
premagic matrices.

\section{Formal certificate and reproducibility}\label{sec:formal}
The accompanying Lean development proves non-diagonalizability using
the standard similarity definition: there do not exist matrices $P,Q$
and a diagonal vector $d$ with
\begin{equation}\label{eq:formal}
 PQ=QP=I,\qquad M=P\diag(d)Q.
\end{equation}
Both inverse identities are stated explicitly. Over a field, this is
equivalent to an invertible change of basis to a diagonal matrix, and
hence to a basis of eigenvectors. The proof first establishes the
kernel-square implication for an arbitrary diagonalizable matrix and
then uses the smaller integer chain
$(1,0,-1)^T\mapsto(3,-3,0)^T\mapsto0$ for $M$. Thus the final Lean
statement certifies the conjecture's actual negation, rather than
assuming a spectral obstruction as a premise.

Further Lean theorems identify this definition with invertible similarity
using the library matrix inverse. They certify all scalar shifts over
characteristic-zero fields and the real normalized stochastic family.
The complex assertion for the normalized family follows from the written
scaling proof above. The characteristic polynomials, determinant formulas,
dimension-minimality observation and arbitrary-dimension extension have
written proofs here; they are not separately formalized in the supplied
Lean file.

The package pins Lean~4.34.1 and the following mathlib commit:
\begin{center}
\texttt{d13f23b723b8a846827a245b89c10fc7d3f11612}.
\end{center}
The verification script records theorem types, their axiom dependencies,
compiler output and file hashes. No \texttt{sorry}, \texttt{admit},
custom axiom, native-decide axiom or stipulated OpenAI theorem is used.
The file \texttt{COVERAGE.md} identifies the certified declarations and
distinguishes them from the explanatory corollaries in this note.
The finite matrix identities, determinant and displayed Jordan form
are also checked independently over the rational numbers.

\section{Attribution and novelty limits}
The source conjecture and definition were checked in the pinned
arXiv version~2 of Teknomo's paper, dated May~6, 2018,
Definitions~4.3--4.5 and Conjecture~7.2~\cite{Teknomo}.
The same conjecture appears in the 2019 journal version~\cite{Journal}.
MathDB problem~336488 supplies the problem-selection record~\cite{MathDB}.
The accompanying provenance record documents the retrievals and the
limited search for prior counterexamples.

The matrix construction uses elementary linear algebra and is not
presented as a new general spectral technique. Absence of a located
earlier refutation does not establish historical priority. The precise
deliverable here is an explicit positive counterexample, its full
written proof, and a kernel-checked formal refutation. The proof is
unconditional and independent of the stipulated corpus used by other
projects in this repository.

\begin{thebibliography}{9}
\bibitem{Teknomo}
K.~Teknomo, \emph{Premagic and Ideal Flow Matrices},
arXiv:1706.08856v2, May 6, 2018.
\url{https://arxiv.org/abs/1706.08856v2}.
Definitions~4.3--4.5; Conjecture~7.2.

\bibitem{Journal}
K.~Teknomo, \emph{Premagic and Ideal Flow Matrices},
Journal of Computing and Applied Informatics \textbf{3}(1) (2019), 35--45.
\url{https://doi.org/10.32734/jocai.v3.i1-621}.
Conjecture~7.2, p.~44.

\bibitem{MathDB}
MathDB, \emph{Diagonalizability conjecture for premagic matrices},
problem 336488. Accessed October 25, 2026.
\url{https://mathdb.com/p/336488}.
\end{thebibliography}
\end{document}
